% Generated draft: requires mathematical and source review.
\documentclass[11pt]{article}
\usepackage[margin=25mm]{geometry}
\usepackage{fontspec}
\setmainfont{Latin Modern Roman}
\usepackage{amsmath,amssymb,mathtools,unicode-math}
\setmathfont{Latin Modern Math}
\usepackage{enumitem,needspace,xurl,hyperref}
\setlength{\parindent}{0pt}
\setlength{\parskip}{5pt}
\setlength{\emergencystretch}{3em}
\newcommand{\N}{\mathbb{N}}
\newcommand{\Z}{\mathbb{Z}}
\newcommand{\Q}{\mathbb{Q}}
\newcommand{\R}{\mathbb{R}}
\newcommand{\C}{\mathbb{C}}
\newcommand{\F}{\mathbb{F}}
\newcommand{\A}{\mathbb{A}}
\newcommand{\PP}{\mathbb{P}}
\newcommand{\E}{\mathbb{E}}
\newcommand{\eps}{\varepsilon}
\newcommand{\dd}{\,\mathrm{d}}
\begin{document}
\section*{UnsolvedMath ID: 6700030}
\textbf{Source problem code:} catalogue entry
\subsection{Definitions and mathematical statement}
The source proposes a simultaneous realization problem for two metric-dependent invariants, K-area and slice-area, on compact simply connected manifolds in every dimension \(n\ge4\). Write these provisionally as \(\operatorname{Karea}(M,g)\) and \(\operatorname{SliceArea}(M,g)\), retaining the source's distinct invariant names.

On the usual positive-value reading of arbitrarily prescribed finite values, the conjectural target is
\[
\forall n\ge4\ \forall a,b\in(0,\infty)\ \exists(M^n,g),\qquad
M\text{ compact and simply connected},
\]
\[
\operatorname{Karea}(M,g)=a,\qquad
\operatorname{SliceArea}(M,g)=b.
\]
Both prescribed values must be realized by the same metric on the same manifold. The manifold and metric may depend on \(n,a,b\); no single fixed topology realizing every pair is asserted. Exact simultaneous prescription is stronger than making both invariants large or approximating a target pair. The phrase probably exist is retained as a conjectural proposal, not a supplied construction.

The record does not define either invariant or its normalization, including a K-area convention in odd dimensions and the precise slicing class in slice-area. It also does not say whether zero is an admissible prescribed value or whether compact manifolds must have empty boundary. The positive-value display makes these unresolved domain choices explicit rather than asserting values outside the natural nonnegative range of area-type invariants. No generic bundle-curvature or sweepout formula is substituted for the missing source definitions.

The supplied research summary calls independent exact prescription open as posed and reports no dedicated verified construction. It provides no additional theorem or reference resolving the stated quantifiers.
\subsection{Short English statement}
In every dimension \(n\ge4\), can one compact simply connected Riemannian manifold realize any prescribed finite K-area and slice-area simultaneously? The invariant definitions, admissible values, and boundary convention are missing.
\subsection{Sources}
\begin{itemize}
\item Local UnsolvedMath registry, permanent ID \(6700030\), exact source code \texttt{catalogue entry}; printed scalar-curvature tag [?30], recorded occurrence 31. Proposer Misha Gromov, 2017.
\item Gromov, \emph{101 Questions, Problems and Conjectures around Scalar Curvature} (2017), recorded PDF page 41; \url{https://www.ihes.fr/~gromov/wp-content/uploads/2018/08/101-problemsOct1-2017.pdf}. Recorded access 2026-07-29; pdf-text-normalized extraction.
\item Supplied OPEN-TRIAGE research summary reports no verified dedicated construction and supplies no additional publication citation or invariant definition.
\end{itemize}
\end{document}
